\chapter{Introduction}

This game is called Deliveroo. Agents move on a tile grid, pick up parcels across the map,
and carry them to delivery tiles to score points.
Parcel rewards decay over time, so the goal is to maximize total delivered score by collecting
and dropping off parcels efficiently before they lose value. There is also the possibility to gain extra points 
by answering questions and following instructions received in the chat by the Admin during the game.
The agents of each team cooperate to compete against the other teams.

The aim is to develop two cooperating agents:
\begin{itemize}
\item Agent 1 — BDI physical executor: does the actual moving/picking/delivering.
\item Agent 2 — LLM cognitive coordinator: has the same function of Agent 1 but extends them with the capability
of interpreting Admin instructions received as messages in the chat, to derive policy rules, coordinating the pair
over P2P and answer Q\&A questions to gather extra points with respect to the simple pick up and deliver process.
\end{itemize}
Other than that PDDL planning is used as an added tool to solve complex problems/situations, such as maps with crates.

All the code and the instuctions to run the agents can be found in the GitHub repository:\\ \url{https://github.com/ale-bena/asa-autobots.git}

\section{Initial state}

Both agents connect to the socket via DjsConnect() and they receive the initial state of the game.
This last consists in the static grid map, containing width, height and the tiles, which are parsed via an enum (\textit{pseudocode}):

\noindent
\begin{minipage}[t]{0.5\textwidth}
\begin{itemize}
    \item WALL
    \item SPAWN
    \item DELIVERY
\end{itemize}
\end{minipage}%
\begin{minipage}[t]{0.5\textwidth}
\begin{itemize}
    \item PAVEMENT
    \item CRATE\_SPAWN, CRATE
    \item ARROW\_UP, ARROW\_RIGHT, ARROW\_DOWN, ARROW\_LEFT
\end{itemize}
\end{minipage}

They also recieve their own identity and state (id, name, x, y, score) \uline{via onYou} and are stored in beliefs.me.
The game parameters, such as observation distance, player capacity, parcel decay interval,
movement duration are received separately via onConfig and are kept in beliefs.config and used for calibration.
The most used ones are cached as dedicated fields.
\uline{Both are part of the same belief base, populated at different points,
and serve to calibrate the agent's reasoning.}
